%%% FIELD proof-prop-honest-scc-and-effect-unions
\textbf{Proof.}
For either split, let \(q\in\mathbb N\), let \(\gamma\in(0,1)\), and write
\[
 S_q=\frac1q\sum_{r=1}^q X_rX_r^\top,
 \qquad X_1,\ldots,X_q\stackrel{\mathrm{iid}}{\sim}\mathcal N(0,\Sigma_0).
\]
We first derive a covariance concentration inequality uniformly over the stated model family.  For indices \(a,b,c,d\in V\), the Gaussian moment-generating function is
\[
 M(t)=\mathbb E_{\Sigma_0}\!\left[\exp(t^\top X_1)\right]
     =\exp\!\left(\frac12t^\top\Sigma_0t\right).
\]
Differentiating this identity four times and evaluating at \(t=0\) gives
\[
 \mathbb E_{\Sigma_0}[X_{1a}X_{1b}X_{1c}X_{1d}]
 =\Sigma_{0,ab}\Sigma_{0,cd}
  +\Sigma_{0,ac}\Sigma_{0,bd}
  +\Sigma_{0,ad}\Sigma_{0,bc}.
\]
Taking \(c=a\) and \(d=b\), and subtracting
\(\mathbb E[X_{1a}X_{1b}]^2=\Sigma_{0,ab}^2\), yields
\[
 \operatorname{Var}_{\Sigma_0}(X_{1a}X_{1b})
 =\Sigma_{0,aa}\Sigma_{0,bb}+\Sigma_{0,ab}^{,2}.
\]
Independence across observations and entrywise unbiasedness of \(S_q\) therefore imply
\[
\begin{aligned}
 \mathbb E_{\Sigma_0}\|S_q-\Sigma_0\|_F^2
 &=\sum_{a,b\in V}\operatorname{Var}_{\Sigma_0}((S_q)_{ab})\\
 &=\frac1q\sum_{a,b\in V}
   \left(\Sigma_{0,aa}\Sigma_{0,bb}+\Sigma_{0,ab}^{,2}\right)\\
 &=\frac{(\operatorname{tr}\Sigma_0)^2+\operatorname{tr}(\Sigma_0^2)}q.
\end{aligned}
\tag{1}
\]
By the covariance-eigenvalue-margin assumption, every eigenvalue of the positive-definite matrix \(\Sigma_0\) is at most \(\overline M\).  Hence
\[
 \operatorname{tr}\Sigma_0\le p\overline M,
 \qquad
 \operatorname{tr}(\Sigma_0^2)\le p\overline M^2.
\]
Substitution into \((1)\) gives
\[
 \mathbb E_{\Sigma_0}\|S_q-\Sigma_0\|_F^2
 \le \frac{p(p+1)\overline M^2}{q}.
\tag{2}
\]
For any nonnegative random variable \(W\) and \(a>0\), the pointwise inequality
\(W\ge a\mathbf 1\{W>a\}\) implies
\(\Pr\{W>a\}\le \mathbb E[W]/a\).  Applying this fact to
\(W=\|S_q-\Sigma_0\|_F^2\), using \((2)\), and invoking the definition
\[
 r_q(\gamma)^2=\frac{p(p+1)\overline M^2}{q\gamma},
\]
we obtain
\[
 \Pr_{\Sigma_0}\!\left\{\|S_q-\Sigma_0\|_F>r_q(\gamma)\right\}
 \le
 \frac{p(p+1)\overline M^2/q}{r_q(\gamma)^2}
 =\gamma.
\tag{3}
\]
The right-hand side of \((3)\) does not depend on \(G_0\) or \(\Sigma_0\); thus the bound is uniform over all \(G_0\in\mathcal F\) and \(\Sigma_0\in\mathcal M_{G_0}\).

Apply \((3)\) to the first split with \((q,\gamma)=(n_1,\alpha_1)\), and define
\[
 \mathcal E_1
 :=\left\{\|S_{n_1}-\Sigma_0\|_F\le r_{n_1}(\alpha_1)\right\}.
\]
The iid Gaussian sampling assumption gives
\(\Pr_{\Sigma_0}(\mathcal E_1)\ge1-\alpha_1\).  Since
\(\Sigma_0\in\mathcal M_{G_0}\), on \(\mathcal E_1\) the definition of distance to a set gives
\[
\begin{aligned}
 \operatorname{dist}_F(S_{n_1},\mathcal M_{G_0})
 &=\inf_{A\in\mathcal M_{G_0}}\|S_{n_1}-A\|_F\\
 &\le \|S_{n_1}-\Sigma_0\|_F\\
 &\le r_{n_1}(\alpha_1).
\end{aligned}
\]
By the definition of \(\widehat{\mathcal G}_{\alpha_1}\), this proves the event inclusion
\[
 \mathcal E_1\subseteq\{G_0\in\widehat{\mathcal G}_{\alpha_1}\},
\]
and consequently
\[
 \Pr_{\Sigma_0}\{G_0\in\widehat{\mathcal G}_{\alpha_1}\}
 \ge1-\alpha_1.
\tag{4}
\]
Whenever \(G_0\in\widehat{\mathcal G}_{\alpha_1}\), its partition \(\operatorname{SCC}(G_0)\) belongs to
\(\{\operatorname{SCC}(G):G\in\widehat{\mathcal G}_{\alpha_1}\}\).  Therefore the event in \((4)\) is contained in the corresponding SCC-coverage event.  Taking the infimum of both probability bounds over \(G_0\in\mathcal F\) and \(\Sigma_0\in\mathcal M_{G_0}\) proves the two asserted first-split inequalities.

We next verify that the effect-set construction is well defined under the stated stability and compactness conditions.  Fix \(G\in\mathcal F\), \((\Lambda_G,\omega)\in\Theta_G\), and distinct \(u,v\in V\).  Let \(J_u:\mathbb R^{V\setminus\{u\}}\to\mathbb R^V\) be the coordinate isometry.  Then
\[
 \Lambda_{G,-u,-u}=J_u^\top\Lambda_GJ_u,
\]
so the stability-margin assumption implies
\[
 \|\Lambda_{G,-u,-u}\|_{\mathrm{op}}
 \le\|\Lambda_G\|_{\mathrm{op}}
 \le1-\kappa<1.
\tag{5}
\]
For \(B=\Lambda_{G,-u,-u}^\top\), the finite geometric identity and \((5)\) give
\[
 (I-B)\sum_{k=0}^N B^k=I-B^{N+1},
 \qquad
 \|B^{N+1}\|_{\mathrm{op}}le(1-\kappa)^{N+1}\longrightarrow0.
\]
Taking the operator-norm limit proves
\[
 (I-\Lambda_{G,-u,-u}^\top)^{-1}
 =\sum_{k=0}^{\infty}(\Lambda_{G,-u,-u}^\top)^k.
\tag{6}
\]
Thus \(\tau_{u\to v}^G(\Lambda_G)\) is finite and continuous in \(\Lambda_G\).  The same argument applied to \(\Lambda_G\) itself shows that \(I_p-\Lambda_G\) is invertible.  Because \(\omega>0\), the covariance map
\[
 (\Lambda_G,\omega)\longmapsto
 \omega(I_p-\Lambda_G)^{-\top}(I_p-\Lambda_G)^{-1}
\]
is well defined, continuous, and positive definite on \(\Theta_G\).

For a realized second-split covariance \(S_{n_2}\), the set \(C_G\) is the intersection of compact \(\Theta_G\), supplied by the compact-parameter assumption, with the inverse image of a closed Frobenius ball under the continuous covariance map.  Hence \(C_G\) is compact.  Whenever \(C_G\ne\varnothing\), continuity from \((6)\) implies that the infimum and supremum defining the \(G\)-indexed interval in \(\mathcal I_{u\to v}\) are finite and attained.  These random endpoints are measurable.  Indeed, for any \(c\in\mathbb R\), the set of symmetric matrices \(S\) for which there exists \((\Lambda,\omega')\in\Theta_G\) satisfying
\[
 \|S-\Sigma_G(\Lambda,\omega')\|_F\le r_{n_2}(\alpha_2)
 \quad\text{and}\quad
 \tau_{u\to v}^G(\Lambda)\le c
\]
is closed: from any convergent sequence of such \(S\)'s, compactness of \(\Theta_G\) supplies a convergent subsequence of witnesses, and continuity preserves both displayed inequalities in the limit.  The same argument applies with \(\tau_{u\to v}^G(\Lambda)\ge c\).  Since \(\mathcal F\) is finite, the unions involved in \(\mathcal I_{u\to v}\) are Borel measurable.

Now fix a true parameter \(\theta_0=(\Lambda_{G_0},\omega)\in\Theta_{G_0}\).  By the meaning of truth in the homoscedastic covariance model, and by the covariance-map and model-image definitions,
\[
 \Sigma_0
 =\Sigma_{G_0}(\Lambda_{G_0},\omega)
 =\omega(I_p-\Lambda_{G_0})^{-\top}(I_p-\Lambda_{G_0})^{-1}
 \in\mathcal M_{G_0}.
\tag{7}
\]
Apply \((3)\) to the second split with \((q,\gamma)=(n_2,\alpha_2)\), and let
\[
 \mathcal E_2
 :=\left\{\|S_{n_2}-\Sigma_0\|_F\le r_{n_2}(\alpha_2)\right\}.
\]
The iid Gaussian sampling assumption gives
\(\Pr_{\Sigma_0}(\mathcal E_2)\ge1-\alpha_2\).  Equation \((7)\) and the definition of \(C_{G_0}\) imply
\[
 \mathcal E_2
 \subseteq
 \{(\Lambda_{G_0},\omega)\in C_{G_0}\}.
\]
It follows that
\[
 \Pr_{\Sigma_0}\{(\Lambda_{G_0},\omega)\in C_{G_0}\}
 \ge1-\alpha_2.
\tag{8}
\]

On \(\mathcal E_1\cap\mathcal E_2\), relation \((4)\) gives
\(G_0\in\widehat{\mathcal G}_{\alpha_1}\), while \((7)\) and the definition of \(\mathcal E_2\) give
\((\Lambda_{G_0},\omega)\in C_{G_0}\), so in particular \(C_{G_0}\ne\varnothing\).  Therefore
\[
 \inf_{(\Lambda,\omega')\in C_{G_0}}
 \tau_{u\to v}^{G_0}(\Lambda)
 \le
 \tau_{u\to v}^{G_0}(\Lambda_{G_0})
 \le
 \sup_{(\Lambda,\omega')\in C_{G_0}}
 \tau_{u\to v}^{G_0}(\Lambda).
\tag{9}
\]
The interval in \((9)\) is one of the intervals in the union defining \(\mathcal I_{u\to v}\).  Hence
\[
 \mathcal E_1\cap\mathcal E_2
 \subseteq
 \{\tau_{u\to v}^{G_0}(\Lambda_{G_0})\in\mathcal I_{u\to v}\}.
\tag{10}
\]
Finally, the union bound, the two instances of \((3)\), and the split-level assumption yield
\[
\begin{aligned}
 \Pr_{\Sigma_0}\{\tau_{u\to v}^{G_0}(\Lambda_{G_0})\in\mathcal I_{u\to v}\}
 &\ge \Pr_{\Sigma_0}(\mathcal E_1\cap\mathcal E_2)\\
 &\ge1-\Pr_{\Sigma_0}(\mathcal E_1^c)-\Pr_{\Sigma_0}(\mathcal E_2^c)\\
 &\ge1-\alpha_1-\alpha_2\\
 &\ge1-\alpha.
\end{aligned}
\]
Independence of the splits ensures the asserted iid Gaussian law separately on each split; the final union-bound step itself needs only the two marginal probability bounds.  Together with \((4)\) and \((8)\), this proves every claim.
